\chapter*{Abstract}

Road-event detection requires identifying agents, their actions, their scene locations, and the compositions formed by these labels. In ROAD-Waymo, unequal training support makes rare compositions difficult to recognize even when overall detection improves. This thesis investigates how decoupled localization, learned composition, and language-informed contextual representations affect recognition of these events.

A YOLOv8x detector supplies candidate boxes and agent identities, while separately trained event heads use frozen video features. The study progresses from overlap-based score transfer and direct region classification to crop-based semantic classification, scene and language attention, and complementary expert blending. The final Contextual Semantic Expert Blend combines a dynamic-class-balancing classification expert with a second expert trained using an additional visual--phrase contrastive objective. A held-out training partition selects one fixed probability-blending weight per seed. Larger paired video and text encoders supply the final system's frozen representations.

Evaluation uses frame-level mean average precision at intersection-over-union 0.5 on 36,717 shared ROAD-Waymo validation frames. The final three-seed mean reaches 15.75\% triplet mAP and 9.00\% mean AP over the 47 triplet classes below the geometric mean training frequency, compared with 7.54\% and 4.14\% for the replicated joint baseline. Relative to the original-encoder DCB blend, triplet mAP increases by 2.32 percentage points; approximately 65\% of that increase comes from tail classes, while common-class performance also improves. Controlled comparisons show that the benefit depends on how language, context, and the loss are combined: phrase-score fusion alone gives small changes, and DCB with the original encoders improves overall performance while reducing tail performance. These results support complementary contextual representations and expert predictions under the studied protocol, with limitations from validation reuse, fixed-vocabulary evaluation, and system-level changes in the scaling comparison.
